% roteiro_entrevistas.tex — fase qualitativa da pesquisa
% Gerado por selecao_entrevistas.py. Compilar com pdflatex (TeX Live 2023).
% Substituições padrão do projeto: sem babel português e sem lmodern;
% \sloppy, \emergencystretch=3em, microtype sem expansão nem protrusão,
% xcolor carregado antes de hyperref.
\documentclass[12pt,a4paper]{article}
\usepackage[T1]{fontenc}
\usepackage[utf8]{inputenc}
\usepackage[top=2.2cm, bottom=2.2cm, left=2.4cm, right=2.4cm]{geometry}
\usepackage{booktabs}
\usepackage{enumitem}
\usepackage[expansion=false,protrusion=false]{microtype}
\usepackage{xcolor}
\usepackage[hidelinks]{hyperref}
\sloppy
\emergencystretch=3em

% linha que liga cada pergunta ao achado quantitativo que ela testa
\newcommand{\achado}[1]{%
  \noindent{\small\itshape\textcolor{black!60}{Achado que a pergunta
  testa: #1}}\par\smallskip}

\title{Roteiro de Entrevistas Semiestruturadas\\
  {\large Fase Qualitativa --- Gest\~ao Hospitalar no SUS/SP, 2015--2025}}
\author{Jo\~ao Eduardo Pastori Garcia (Estat\'istico Respons\'avel)}
\date{Julho de 2026}

\begin{document}
\maketitle
\tableofcontents
\clearpage

\section{Apresentação e método}

Este roteiro operacionaliza a fase qualitativa da pesquisa sobre o efeito
do modelo institucional de gestão no desempenho de hospitais SUS de São
Paulo (painel balanceado de 289 hospitais, 2015 a 2025). A fase de
estimação, concluída, deixou quatro achados que as entrevistas devem
explicar por dentro: (i) o regime operacional mais intenso e a eficiência
técnica maior das OSS (DEA-BCC corrigida de viés e SFA-BC95) são traços
robustos da rede; (ii) o tempo médio de permanência cerca de 10\% menor
nas OSS é evidência sugestiva de variação dentro do hospital, apoiada em
apenas 5 conversores; (iii) a mortalidade menor concentra-se na
trajetória dos conversores e não sobrevive ao controle sintético, o que
recomenda cautela e investigação de mecanismo; (iv) o faturamento real por
saída não difere entre categorias --- a diferença de gestão, se existe,
não aparece no faturamento. Ressalva de leitura (regime de pagamento):
hospitais OSS são custeados por contrato de gestão com orçamento global,
não por reembolso de AIH --- o valor registrado por saída é, para eles,
produção sem consequência financeira direta; a diferença de faturamento
pode refletir o próprio regime de pagamento, a variável cujo efeito se
busca medir, e não uso de recursos ou qualidade assistencial. Toda
leitura de mortalidade nesta fase usa a
complexidade estrutural como referência de perfil, nunca a complexidade
ponderada pela mortalidade, para evitar circularidade.

A seleção de entrevistados segue seis blocos amostrais descritos na
Seção~\ref{sec:selecao}: os 5 conversores Direta$\to$OSS (obrigatórios),
desviantes positivos e negativos de eficiência persistente no suporte
comum (faixas Barcelona 3--4, portes médio e grande), pares casados de
contraste OSS vs.\ Direta e OSS vs.\ Filantrópico com estrutura parecida
e eficiência divergente, casos de maior divergência entre DEA e SFA, e
duas entrevistas curtas de validação de registro. O grupo Privado (3
hospitais) fica fora dos blocos comparativos, pelo tamanho e pelo
perfil assistencial não comparável (o painel de 289 não tem mais
hospitais de longa permanência: a ETAPA F removeu todos os 18 que
existiam no painel de 314). As
entrevistas são semiestruturadas: o roteiro-base da Seção~\ref{sec:base}
vale para todos os entrevistados dos blocos a--e, e os módulos da
Seção~\ref{sec:modulos} acrescentam perguntas específicas por perfil.
Cada pergunta traz, em itálico, o achado quantitativo que ela testa ---
o entrevistador deve conhecer esses achados, mas não deve revelá-los ao
entrevistado antes da resposta.

\section{Hospitais selecionados}
\label{sec:selecao}

\subsection{Bloco a --- Conversores Direta$\to$OSS (obrigatórios)}
\begin{center}\small
\begin{tabular}{llllrr}
\toprule
CNES & Hospital & Município & Categoria & DEA & SFA \\
\midrule
2081695 & Conjunto Hospitalar Sorocaba & Sorocaba & OSS (conversor) & 0,563 & 0,999 \\
2078287 & Centro De Referencia Da Saude Da Mulhe & Sao Paulo & OSS (conversor) & 0,827 & 0,805 \\
2082225 & Hospital Katia De Souza Rodrigues Taip & Sao Paulo & OSS (conversor) & 0,818 & 0,912 \\
2091755 & Hospital Geral De Vila Penteado Dr Jos & Sao Paulo & OSS (conversor) & 0,582 & 0,893 \\
2750511 & Hospital Estadual Dr Odilo Antunes De  & Presidente Prudente & OSS (conversor) & 0,698 & 0,900 \\
\bottomrule
\end{tabular}
\end{center}
\noindent\textbf{2081695 --- Conjunto Hospitalar Sorocaba} (Sorocaba, interior; Grande Porte, faixa Barcelona 5, complexidade estrutural 675,0). Conversor Direta->OSS em 2019. Antes/depois (medianas): TMP (dias): 8,29 -> 7,13 (-14,0\%); Mortalidade geral: 0,09 -> 0,08 (-14,8\%); Faturamento real por saída (R\$ de 2025): 2.734,63 -> 2.401,47 (-12,2\%); Saídas hospitalares (produção): 13.159,50 -> 16.243,00 (23,4\%)\par\medskip
\noindent\textbf{2078287 --- Centro De Referencia Da Saude Da Mulher Sao Paulo} (Sao Paulo, capital; Médio Porte, faixa Barcelona 4, complexidade estrutural 266,0). Conversor Direta->OSS em 2023. Antes/depois (medianas): TMP (dias): 4,12 -> 3,40 (-17,5\%); Mortalidade geral: 0,08 -> 0,05 (-39,5\%); Faturamento real por saída (R\$ de 2025): 1.989,40 -> 1.423,76 (-28,4\%); Saídas hospitalares (produção): 5.583,00 -> 14.036,00 (151,4\%)\par\medskip
\noindent\textbf{2082225 --- Hospital Katia De Souza Rodrigues Taipas Sao Paulo} (Sao Paulo, capital; Grande Porte, faixa Barcelona 4, complexidade estrutural 380,0). Conversor Direta->OSS em 2025 [apenas 1 ano pós: sem leitura de tendência]. Antes/depois (medianas): TMP (dias): 5,41 -> 5,35 (-1,1\%); Mortalidade geral: 0,08 -> 0,06 (-18,1\%); Faturamento real por saída (R\$ de 2025): 869,18 -> 977,07 (12,4\%); Saídas hospitalares (produção): 9.804,00 -> 10.839,00 (10,6\%)\par\medskip
\noindent\textbf{2091755 --- Hospital Geral De Vila Penteado Dr Jose Pangella Sao Paulo} (Sao Paulo, capital; Grande Porte, faixa Barcelona 4, complexidade estrutural 350,0). Conversor Direta->OSS em 2025 [apenas 1 ano pós: sem leitura de tendência]. Antes/depois (medianas): TMP (dias): 5,03 -> 5,50 (9,3\%); Mortalidade geral: 0,07 -> 0,08 (11,1\%); Faturamento real por saída (R\$ de 2025): 940,08 -> 1.246,13 (32,6\%); Saídas hospitalares (produção): 8.350,50 -> 7.783,00 (-6,8\%)\par\medskip
\noindent\textbf{2750511 --- Hospital Estadual Dr Odilo Antunes De Siqueira P Prudente} (Presidente Prudente, interior; Médio Porte, faixa Barcelona 2, complexidade estrutural 94,0). Conversor Direta->OSS em 2025 [apenas 1 ano pós: sem leitura de tendência]. Antes/depois (medianas): TMP (dias): 3,72 -> 3,91 (5,4\%); Mortalidade geral: 0,01 -> 0,01 (13,4\%); Faturamento real por saída (R\$ de 2025): 1.132,81 -> 1.082,27 (-4,5\%); Saídas hospitalares (produção): 3.675,00 -> 3.513,00 (-4,4\%)\par\medskip

\subsection{Bloco b --- Desviantes positivos (quartil superior persistente)}
\begin{center}\small
\begin{tabular}{llllrr}
\toprule
CNES & Hospital & Município & Categoria & DEA & SFA \\
\midrule
7573162 & Hospital Regional De Jundiai & Jundiai & OSS & 0,905 & 0,999 \\
6164366 & Hospital Estadual Americo Brasiliense & Americo Brasiliense & Direta & 0,810 & 0,908 \\
5718368 & Hosp Mun M Boi Mirim & Sao Paulo & Público Municipal & 0,870 & 0,962 \\
2076896 & Hospital Sao Luiz Gonzaga & Sao Paulo & Filantrópico & 0,922 & 0,928 \\
\bottomrule
\end{tabular}
\end{center}
\noindent\textbf{7573162 --- Hospital Regional De Jundiai} (Jundiai, interior; Médio Porte, faixa Barcelona 3, complexidade estrutural 246,0). Desviante positivo da categoria OSS (quartil superior persistente de te\_bcc\_vc e te\_sfa): eficiência DEA 0,905 e SFA 0,999 (medianas 2015-2025), no suporte comum (faixa 3, Médio Porte). Marcador financeiro x físico: 4º no rank DEA com insumo financeiro contra 18º só com leitos (diferença de 14 posições; positivo = ranqueia melhor com o insumo financeiro do que sem ele) — investigar se a divergência é regime de pagamento, porte, complexidade ou gestão real.\par\medskip
\noindent\textbf{6164366 --- Hospital Estadual Americo Brasiliense} (Americo Brasiliense, interior; Médio Porte, faixa Barcelona 3, complexidade estrutural 187,0). Desviante positivo da categoria Direta (critério relaxado: quartil superior de DEA e acima da mediana de SFA): eficiência DEA 0,810 e SFA 0,908 (medianas 2015-2025), no suporte comum (faixa 3, Médio Porte). Marcador financeiro x físico: 27º no rank DEA com insumo financeiro contra 75º só com leitos (diferença de 48 posições; positivo = ranqueia melhor com o insumo financeiro do que sem ele) — investigar se a divergência é regime de pagamento, porte, complexidade ou gestão real.\par\medskip
\noindent\textbf{5718368 --- Hosp Mun M Boi Mirim} (Sao Paulo, capital; Grande Porte, faixa Barcelona 4, complexidade estrutural 469,0). Desviante positivo da categoria Público Municipal (quartil superior persistente de te\_bcc\_vc e te\_sfa): eficiência DEA 0,870 e SFA 0,962 (medianas 2015-2025), no suporte comum (faixa 4, Grande Porte). Marcador financeiro x físico: 12º no rank DEA com insumo financeiro contra 4º só com leitos (diferença de -8 posições; positivo = ranqueia melhor com o insumo financeiro do que sem ele) — investigar se a divergência é regime de pagamento, porte, complexidade ou gestão real.\par\medskip
\noindent\textbf{2076896 --- Hospital Sao Luiz Gonzaga} (Sao Paulo, capital; Grande Porte, faixa Barcelona 3, complexidade estrutural 241,0). Desviante positivo da categoria Filantrópico (quartil superior persistente de te\_bcc\_vc e te\_sfa): eficiência DEA 0,922 e SFA 0,928 (medianas 2015-2025), no suporte comum (faixa 3, Grande Porte). Marcador financeiro x físico: 1º no rank DEA com insumo financeiro contra 66º só com leitos (diferença de 65 posições; positivo = ranqueia melhor com o insumo financeiro do que sem ele) — investigar se a divergência é regime de pagamento, porte, complexidade ou gestão real.\par\medskip

\subsection{Bloco c --- Desviantes negativos (quartil inferior persistente)}
\begin{center}\small
\begin{tabular}{llllrr}
\toprule
CNES & Hospital & Município & Categoria & DEA & SFA \\
\midrule
2028840 & Instituto De Infectologia Emilio Ribas & Sao Paulo & Direta & 0,447 & 0,409 \\
2698471 & Secao Hospital Municipal Dr Arthur Dom & Santos & Público Municipal & 0,378 & 0,717 \\
2080354 & Hospital Santo Antonio Santos & Santos & Filantrópico & 0,351 & 0,399 \\
\bottomrule
\end{tabular}
\end{center}
\noindent\textbf{2028840 --- Instituto De Infectologia Emilio Ribas Sao Paulo} (Sao Paulo, capital; Grande Porte, faixa Barcelona 4, complexidade estrutural 459,0). Desviante negativo da categoria Direta (quartil inferior persistente de te\_bcc\_vc e te\_sfa): eficiência DEA 0,447 e SFA 0,409 (medianas 2015-2025), no suporte comum (faixa 4, Grande Porte).\par\medskip
\noindent\textbf{2698471 --- Secao Hospital Municipal Dr Arthur Domingues Pinto} (Santos, interior; Médio Porte, faixa Barcelona 3, complexidade estrutural 122,0). Desviante negativo da categoria Público Municipal (quartil inferior persistente de te\_bcc\_vc e te\_sfa): eficiência DEA 0,378 e SFA 0,717 (medianas 2015-2025), no suporte comum (faixa 3, Médio Porte).\par\medskip
\noindent\textbf{2080354 --- Hospital Santo Antonio Santos} (Santos, interior; Grande Porte, faixa Barcelona 3, complexidade estrutural 126,0). Desviante negativo da categoria Filantrópico (quartil inferior persistente de te\_bcc\_vc e te\_sfa): eficiência DEA 0,351 e SFA 0,399 (medianas 2015-2025), no suporte comum (faixa 3, Grande Porte).\par\medskip

\subsection{Bloco d --- Pares casados de contraste}
\begin{center}\small
\begin{tabular}{llllrr}
\toprule
CNES & Hospital & Município & Categoria & DEA & SFA \\
\midrule
2065665 & Hospital Maternidade Interlagos & Sao Paulo & Direta & 0,690 & 0,900 \\
2079232 & Hospital Santa Barbara & Santa Barbara D'Oeste & Filantrópico & 0,523 & 0,820 \\
2792141 & Hospital Regional De Cotia & Cotia & OSS & 0,893 & 0,999 \\
2079410 & Complexo Hospitalar Padre Bento De Gua & Guarulhos & Direta & 0,527 & 0,815 \\
2096412 & Santa Casa De Misericordia De Jacarei & Jacarei & Filantrópico & 0,598 & 0,890 \\
\bottomrule
\end{tabular}
\end{center}
\noindent\textbf{2065665 --- Hospital Maternidade Interlagos} (Sao Paulo, capital; Médio Porte, faixa Barcelona 3, complexidade estrutural 160,0). Par casado da âncora OSS Hospital Regional De Jundiai (CNES 7573162): mesmo porte (Médio Porte), faixa Barcelona 3 vs 3, leitos 89 vs 136, mas eficiência divergente: DEA 0,690 vs 0,905; SFA 0,900 vs 0,999. Estrutura parecida, gestão Direta: o que a gestão faz de diferente?\par\medskip
\noindent\textbf{2079232 --- Hospital Santa Barbara} (Santa Barbara D'Oeste, interior; Médio Porte, faixa Barcelona 3, complexidade estrutural 247,0). Par casado da âncora OSS Hospital Regional De Jundiai (CNES 7573162): mesmo porte (Médio Porte), faixa Barcelona 3 vs 3, leitos 137 vs 136, mas eficiência divergente: DEA 0,523 vs 0,905; SFA 0,820 vs 0,999. Estrutura parecida, gestão Filantrópico: o que a gestão faz de diferente?\par\medskip
\noindent\textbf{2792141 --- Hospital Regional De Cotia} (Cotia, interior; Médio Porte, faixa Barcelona 4, complexidade estrutural 330,0). Âncora OSS de alta eficiência (DEA 0,893, SFA 0,999) para os pares casados.\par\medskip
\noindent\textbf{2079410 --- Complexo Hospitalar Padre Bento De Guarulhos} (Guarulhos, interior; Médio Porte, faixa Barcelona 4, complexidade estrutural 327,0). Par casado da âncora OSS Hospital Regional De Cotia (CNES 2792141): mesmo porte (Médio Porte), faixa Barcelona 4 vs 4, leitos 145 vs 145, mas eficiência divergente: DEA 0,527 vs 0,893; SFA 0,815 vs 0,999. Estrutura parecida, gestão Direta: o que a gestão faz de diferente?\par\medskip
\noindent\textbf{2096412 --- Santa Casa De Misericordia De Jacarei} (Jacarei, interior; Médio Porte, faixa Barcelona 4, complexidade estrutural 312,0). Par casado da âncora OSS Hospital Regional De Cotia (CNES 2792141): mesmo porte (Médio Porte), faixa Barcelona 4 vs 4, leitos 143 vs 145, mas eficiência divergente: DEA 0,598 vs 0,893; SFA 0,890 vs 0,999. Estrutura parecida, gestão Filantrópico: o que a gestão faz de diferente?\par\medskip

\subsection{Bloco e --- Divergência DEA$\times$SFA}
\begin{center}\small
\begin{tabular}{llllrr}
\toprule
CNES & Hospital & Município & Categoria & DEA & SFA \\
\midrule
2077590 & Ibcc & Sao Paulo & Filantrópico & 0,748 & 0,456 \\
2089785 & Hospital Do Rim E Hipertensao & Sao Paulo & Filantrópico & 0,782 & 0,675 \\
2083086 & Hospital Amaral Carvalho Jau & Jau & Filantrópico & 0,801 & 0,765 \\
2790734 & Instituto Lauro De Souza Lima Bauru & Bauru & Direta & 0,742 & 0,642 \\
9601 & Hospital Pio Xii & Sao Jose Dos Campos & Filantrópico & 0,729 & 0,671 \\
\bottomrule
\end{tabular}
\end{center}
\noindent\textbf{2077590 --- Ibcc} (Sao Paulo, capital; Grande Porte, faixa Barcelona 3, complexidade estrutural 176,0). Divergência DEA x SFA: rank DEA 57 vs rank SFA 273 entre 280 CNES (|dif| = 216; DEA favorável, SFA desfavorável; te\_bcc\_vc 0,748, te\_sfa 0,456). Os dois métodos discordam: há algo que os modelos não capturam.\par\medskip
\noindent\textbf{2089785 --- Hospital Do Rim E Hipertensao} (Sao Paulo, capital; Grande Porte, faixa Barcelona 3, complexidade estrutural 227,0). Divergência DEA x SFA: rank DEA 39 vs rank SFA 251 entre 280 CNES (|dif| = 212; DEA favorável, SFA desfavorável; te\_bcc\_vc 0,782, te\_sfa 0,675). Os dois métodos discordam: há algo que os modelos não capturam.\par\medskip
\noindent\textbf{2083086 --- Hospital Amaral Carvalho Jau} (Jau, interior; Grande Porte, faixa Barcelona 4, complexidade estrutural 415,0). Divergência DEA x SFA: rank DEA 28 vs rank SFA 232 entre 280 CNES (|dif| = 204; DEA favorável, SFA desfavorável; te\_bcc\_vc 0,801, te\_sfa 0,765). Os dois métodos discordam: há algo que os modelos não capturam.\par\medskip
\noindent\textbf{2790734 --- Instituto Lauro De Souza Lima Bauru} (Bauru, interior; Médio Porte, faixa Barcelona 2, complexidade estrutural 27,0). Divergência DEA x SFA: rank DEA 67 vs rank SFA 262 entre 280 CNES (|dif| = 195; DEA favorável, SFA desfavorável; te\_bcc\_vc 0,742, te\_sfa 0,642). Os dois métodos discordam: há algo que os modelos não capturam.\par\medskip
\noindent\textbf{9601 --- Hospital Pio Xii} (Sao Jose Dos Campos, interior; Médio Porte, faixa Barcelona 3, complexidade estrutural 141,0). Divergência DEA x SFA: rank DEA 73 vs rank SFA 253 entre 280 CNES (|dif| = 180; DEA favorável, SFA desfavorável; te\_bcc\_vc 0,729, te\_sfa 0,671). Os dois métodos discordam: há algo que os modelos não capturam.\par\medskip

\subsection{Bloco f --- Validação de registro (entrevistas curtas)}
\begin{center}\small
\begin{tabular}{llllrr}
\toprule
CNES & Hospital & Município & Categoria & DEA & SFA \\
\midrule
2022648 & Hospital Irmaos Penteado E Santa Casa  & Campinas & Filantrópico & 0,354 & 0,664 \\
2097613 & Hospital Infante D Henrique & Sao Jose Do Rio Preto & Filantrópico & 0,758 & 0,264 \\
\bottomrule
\end{tabular}
\end{center}
\noindent\textbf{2022648 --- Hospital Irmaos Penteado E Santa Casa De Campinas} (Campinas, interior; Médio Porte, faixa Barcelona 3, complexidade estrutural 156,0). ocupação de UTI de 875\% em 2021: verificar leitos CNES cadastrados vs leitos operacionais na pandemia. Ocupação de UTI registrada: 2020 = 561,6\%, 2021 = 875,2\% (leitos UTI cadastrados: 16).\par\medskip
\noindent\textbf{2097613 --- Hospital Infante D Henrique} (Sao Jose Do Rio Preto, interior; Grande Porte, faixa Barcelona 2, complexidade estrutural 39,0). denominadores frágeis em 2020-21 (menos de 30 saídas sem COVID no biênio, TMP zero em 2021): verificar registro de produção no SIH. Saídas registradas: 2020 = 13, 2021 = 4; TMP 2021 = 0,00.\par\medskip

\section{Roteiro-base comum (blocos a--e)}
\label{sec:base}

O roteiro-base cobre seis domínios e dura de 50 a 60 minutos. As perguntas são abertas; os itens entre colchetes são estímulos de aprofundamento para o entrevistador.

\subsection{Governança e desenho institucional}
\begin{enumerate}[itemsep=2pt]
  \item Como se organiza a direção do hospital: quem decide o quê entre a diretoria local, a mantenedora ou OSS e a Secretaria? [organograma real vs.\ formal; conselhos; tempo de mandato da direção]
  \par\nopagebreak\achado{a eficiência técnica maior das OSS é traço robusto da rede (DEA e SFA concordam); a hipótese é que o desenho de governança seja o canal.}
  \item Que decisões precisam de autorização externa ao hospital e quanto tempo levam, do pedido à resposta? [contratação, compra acima de um teto, abertura de leito]
  \par\nopagebreak\achado{o regime operacional mais intenso das OSS sugere menor fricção decisória; a pergunta mede a fricção diretamente.}
  \item Com que frequência a direção examina indicadores assistenciais e o que acontece quando um indicador piora? [rotina de reunião; quem responde pelo indicador]
  \par\nopagebreak\achado{desviantes positivos mantêm eficiência alta por 11 anos; persistência sugere rotina gerencial, não sorte.}
\end{enumerate}

\subsection{Autonomia de recursos humanos e compras}
\begin{enumerate}[itemsep=2pt]
  \item Como o hospital contrata e desliga pessoal assistencial, e quanto tempo passa entre a necessidade identificada e o profissional trabalhando? [concurso vs.\ CLT; terceirização; banco de horas]
  \par\nopagebreak\achado{a produção das OSS é cerca de 45\% maior a estrutura comparável; elasticidade de RH é o mecanismo candidato.}
  \item Como funcionam as compras e a manutenção de equipamentos: prazos típicos, o que mais trava? [licitação vs.\ compra direta; estoque; equipamento parado]
  \par\nopagebreak\achado{o faturamento real por saída não difere entre categorias; se as OSS produzem mais com faturamento igual, a diferença deve estar no uso dos insumos, não no preço.}
  \item Qual é a rotatividade das equipes médica e de enfermagem, e o que a explica? [vínculos múltiplos; plantões; salário relativo]
  \par\nopagebreak\achado{estabilidade de equipe é explicação concorrente à gestão para a eficiência persistente dos desviantes positivos.}
\end{enumerate}

\subsection{Contrato de gestão e metas}
\begin{enumerate}[itemsep=2pt]
  \item Existe contrato, convênio ou instrumento com metas quantitativas? Quais indicadores, quem apura, e o que acontece quando a meta não é cumprida? [glosa; aditivo; renegociação]
  \par\nopagebreak\achado{o contrato de gestão é a diferença institucional formal entre OSS e administração direta; TMP $-$10,2\% e produção $+$43,7\% são os efeitos que ele teria de explicar.}
  \item As metas chegam ao cotidiano clínico? Dê um exemplo de decisão de ponta (alta, agenda cirúrgica, escala) que mudou por causa de meta. [quem traduz a meta para a equipe]
  \par\nopagebreak\achado{o TMP menor nas OSS é evidência within (5 conversores): a pergunta busca o elo micro entre contrato e permanência.}
\end{enumerate}

\subsection{Relação com a Secretaria e financiamento}
\begin{enumerate}[itemsep=2pt]
  \item Como é a interação com a Secretaria (estadual ou municipal): cobrança, apoio técnico, regulação de fila? Com que frequência? [câmaras técnicas; visitas; sistemas]
  \par\nopagebreak\achado{categorias públicas e OSS respondem à mesma Secretaria com resultados diferentes; a pergunta separa o vínculo formal da relação praticada.}
  \item Como o dinheiro chega e o que acontece numa crise de caixa? [teto MAC; incentivos; atraso de repasse; quem socorre]
  \par\nopagebreak\achado{o faturamento real por saída não difere entre categorias: se o financiamento é semelhante, a diferença de desempenho não é explicada por mais dinheiro.}
\end{enumerate}

\subsection{Gestão de leitos e fluxo do paciente}
\begin{enumerate}[itemsep=2pt]
  \item Descreva a rotina de altas: horário, critérios, quem decide, existe reunião diária de leitos? [alta assistida; previsão de alta na admissão]
  \par\nopagebreak\achado{TMP cerca de 10\% menor nas OSS: a rotina de gestão de altas é o mecanismo mais direto que a entrevista pode verificar.}
  \item Quem controla a fila de internação e a ocupação do leito (NIR, regulação própria, CROSS)? Como se decide quem entra? [leito de retaguarda; internação social]
  \par\nopagebreak\achado{o giro alto dos desviantes positivos exige controle fino de entrada e saída; a pergunta identifica o instrumento.}
  \item Que parte da cirurgia migrou para o regime ambulatorial ou hospital-dia, e o que impediu de migrar mais? [centro cirúrgico; anestesia; retorno]
  \par\nopagebreak\achado{produção maior com TMP menor é compatível com substituição de internação por ambulatório; a pergunta testa essa composição.}
\end{enumerate}

\subsection{Resposta à pandemia (2020--2021)}
\begin{enumerate}[itemsep=2pt]
  \item O que mudou em leitos, equipes e fluxo em 2020--21, e quanto da expansão foi física de fato vs.\ apenas cadastro? [leitos de campanha; UTI improvisada; registro no CNES]
  \par\nopagebreak\achado{as séries de 2020--21 têm extremos de ocupação de UTI e denominadores frágeis; a distinção entre leito real e leito cadastrado é decisiva para a leitura desses anos.}
  \item O que a pandemia deixou de permanente na operação do hospital? [telessaúde; protocolos; equipes]
  \par\nopagebreak\achado{a estimação de robustez sem 2020--21 preserva os achados principais; a pergunta explora se houve mudança estrutural que os dados pós-2021 já refletem.}
\end{enumerate}

\section{Módulos específicos por perfil}
\label{sec:modulos}

\subsection{Módulo A --- Conversores Direta$\to$OSS (bloco a)}

Reconstituir a linha do tempo da transição: quando a conversão começou a
ser discutida, quando o contrato foi assinado, quando a OSS assumiu de
fato. Atenção metodológica: no Pérola Byington e em Sorocaba a melhora
dos indicadores começou \emph{antes} da conversão formal --- é
indispensável perguntar o que já estava em curso no período anterior
(intervenção da Secretaria, troca de direção, obra, mudança de perfil),
pois a resposta separa efeito da gestão OSS de tendência pré-existente.

\begin{enumerate}[itemsep=2pt]
  \item Conte a linha do tempo da transição para OSS: primeira discussão, decisão, assinatura, assunção. O que aconteceu no hospital em cada etapa? [greves; judicialização; transição de direção]
  \par\nopagebreak\achado{as datas de conversão da proxy (2019, 2023 e três em 2025) vêm do cadastro; a linha do tempo real pode antecedê-las e reinterpretar o estudo de eventos.}
  \item O que mudou concretamente no primeiro ano de OSS: rotinas, chefias, sistemas, escalas? O que NÃO mudou? [primeiros 90 dias; o que a OSS trouxe pronto]
  \par\nopagebreak\achado{no estudo de eventos, o TMP cai já nos primeiros anos pós; a pergunta identifica qual prática explica a queda rápida.}
  \item Quanto da equipe assistencial e das chefias permaneceu após a conversão? [estatutários cedidos; recontratação CLT]
  \par\nopagebreak\achado{continuidade de equipe distingue efeito de gestão de efeito de composição --- se a equipe é a mesma, a mudança é de método, não de pessoas.}
  \item Quais metas o contrato de gestão fixou no início e como elas mudaram? Alguma meta era vista como inatingível? [metas de produção vs.\ qualidade]
  \par\nopagebreak\achado{produção $+$43,7\% (transversal) e TMP $-$10,2\% within: as metas contratuais são a explicação institucional candidata.}
  \item (Somente Pérola Byington e Sorocaba) Os indicadores já vinham melhorando antes da conversão formal. O que estava acontecendo no hospital nos 2--3 anos anteriores? [gestor interino; plano de saneamento; antecipação da OSS]
  \par\nopagebreak\achado{a melhora pré-conversão observada nesses dois hospitais ameaça a leitura causal do estudo de eventos; a resposta decide entre antecipação, tendência ou seleção.}
\end{enumerate}

\subsection{Módulo B --- Desviantes positivos (bloco b)}

O objetivo é extrair práticas concretas e transferíveis por trás do giro
alto e da eficiência persistente, com ceticismo ativo: pedir sempre o
exemplo, o documento, a rotina --- não aceitar declaração de princípios.

\begin{enumerate}[itemsep=2pt]
  \item O hospital está no quartil superior de eficiência da rede há 11 anos. Na sua leitura, quais três práticas explicam isso? Para cada uma: desde quando, quem implantou, o que custou? [pedir exemplos datados]
  \par\nopagebreak\achado{a persistência no quartil superior em DEA e SFA simultaneamente afasta artefato de método; deve haver prática real por trás.}
  \item Como funciona, passo a passo, a gestão de altas num dia típico? [huddle; previsão de alta; alta aos sábados]
  \par\nopagebreak\achado{giro alto com TMP baixo é a assinatura operacional dos desviantes positivos.}
  \item Que procedimentos o hospital faz em regime ambulatorial que vizinhos comparáveis internam? [cirurgia ambulatorial; hospital-dia; protocolos]
  \par\nopagebreak\achado{a produção alta com leitos comparáveis sugere substituição ambulatorial; a pergunta lista os procedimentos.}
  \item Como o hospital se relaciona com a regulação de vagas: aceita tudo que a central manda? Recusa o quê, e com que argumento? [seleção de casos; perfil de demanda]
  \par\nopagebreak\achado{eficiência alta pode ser gestão ou seleção de pacientes mais leves; a pergunta confronta a hipótese de seleção --- ler junto com a complexidade estrutural do hospital.}
\end{enumerate}

\subsection{Módulo C --- Desviantes negativos (bloco c)}

Simétrico ao módulo B, mas centrado nas restrições percebidas. O tom não
é de auditoria: o objetivo é entender o que impede o giro, na visão de
quem opera, e confrontar depois com as restrições objetivas do painel.

\begin{enumerate}[itemsep=2pt]
  \item Quais são hoje os três maiores obstáculos para o hospital girar mais os leitos? [financiamento; RH; perfil de demanda; física do prédio]
  \par\nopagebreak\achado{o hospital está no quartil inferior persistente de DEA e SFA; a pergunta colhe a explicação nativa antes de confrontá-la com os dados.}
  \item Descreva o paciente típico que fica internado mais tempo do que precisaria. Por que ele não sai? [internação social; retaguarda; exame que não sai]
  \par\nopagebreak\achado{TMP alto com ocupação alta indica leito ocupado por paciente que não precisaria dele; a pergunta identifica o gargalo.}
  \item O financiamento cobre o custo do que o hospital produz? Onde aperta primeiro? [tabela SUS; folha; incentivos]
  \par\nopagebreak\achado{o faturamento real por saída não difere entre categorias na rede; se a restrição percebida é financeira, ela precisa ser conciliada com esse achado.}
  \item Se o hospital pudesse mudar uma única regra (de RH, compra ou contrato), qual mudaria e o que aconteceria? [pergunta projetiva]
  \par\nopagebreak\achado{a resposta revela qual margem de autonomia o gestor acredita que falta --- comparável, em espelho, com o que as OSS dizem ter.}
\end{enumerate}

\subsection{Módulo D --- Pares casados (bloco d)}

As perguntas dos domínios 2 (RH e compras), 3 (contrato e metas) e 5
(leitos e fluxo) do roteiro-base devem ser feitas de forma espelhada e
com a mesma ordem nos dois hospitais do par, para permitir comparação
direta das respostas. Acrescentar:

\begin{enumerate}[itemsep=2pt]
  \item Num dia típico, quem decide a abertura de um leito extra, a contratação de um plantonista e a compra urgente de um insumo --- e em quanto tempo? [responder para os três casos]
  \par\nopagebreak\achado{o par tem faixa Barcelona, porte e leitos parecidos com eficiência divergente; a velocidade decisória é a diferença de gestão mais plausível com estrutura igual.}
  \item Como é montada a escala médica do fim de semana e o que acontece com as altas de sexta a domingo? [alta no fim de semana; plantão de sexta]
  \par\nopagebreak\achado{diferenças de TMP entre pares estruturalmente parecidos costumam se concentrar no fim de semana; pergunta espelhada nos dois hospitais do par.}
\end{enumerate}

\subsection{Módulo E --- Divergência DEA$\times$SFA (bloco e)}

A divergência entre os dois métodos sugere algo que os modelos não
capturam: composição de casos atípica, insumo mal medido, produção
heterogênea. A entrevista busca esse fator omitido.

\begin{enumerate}[itemsep=2pt]
  \item O que este hospital faz que hospitais parecidos não fazem --- serviços, perfis de paciente, papéis na rede que não aparecem em leitos e saídas? [ensino; referência regional; porta aberta vs.\ referenciada]
  \par\nopagebreak\achado{DEA e SFA discordam fortemente sobre este hospital; a explicação provável é dimensão de produto ou insumo que nenhum dos dois modelos mede.}
  \item Os números oficiais de leitos, produção e permanência do hospital refletem a operação real? Onde o registro distorce? [leitos bloqueados; produção não faturada]
  \par\nopagebreak\achado{a divergência entre métodos também pode ser artefato de medida; a pergunta verifica a qualidade do dado antes de interpretar o resto da entrevista.}
\end{enumerate}

\subsection{Módulo F --- Validação de registro (bloco f, entrevistas curtas)}

Entrevistas de 20 a 30 minutos, apenas com o responsável pelo cadastro
CNES e faturamento SIH, restritas ao registro --- não são entrevistas de
gestão e não usam o roteiro-base.

\begin{enumerate}[itemsep=2pt]
  \item (CNES 2022648) Em 2021 a ocupação de UTI calculada chega a 875\%. Quantos leitos de UTI operavam de fato em 2020--21 e quantos estavam no CNES? Houve atraso ou represamento de atualização cadastral? [leitos de campanha; UTI COVID]
  \par\nopagebreak\achado{ocupação acima de 100\% em escala implica denominador (leitos CNES) menor que a operação real; a resposta define como tratar as observações extremas de 2020--21.}
  \item (CNES 2097613) O registro mostra menos de 30 saídas sem COVID no biênio 2020--21 e TMP zero em 2021. O hospital fechou, mudou de papel ou parou de faturar pelo SIH nesse período? [interdição; reforma; migração de convênio]
  \par\nopagebreak\achado{as duas observações desse CNES já são tratadas como denominador frágil e excluídas dos desfechos; a resposta confirma ou corrige esse tratamento.}
\end{enumerate}

\section{Logística e esforço}

A seleção soma 24 hospitais únicos: 22 entrevistas
plenas (roteiro-base de 50--60 minutos mais módulo específico de 20--30
minutos) e 2 entrevistas curtas de validação de registro (20--30
minutos). O teto recomendado é de 15 a 18 entrevistas. Havendo
necessidade de corte, a ordem de prioridade é: bloco a (obrigatório),
bloco f (curtas, custo baixo), blocos b e c (o contraste
positivo--negativo é o núcleo da fase), bloco d (manter ao menos um par
completo por âncora, preferindo o par Direta, que espelha a comparação
principal da estimação) e, por último, bloco e. Recomenda-se gravar com
consentimento, transcrever integralmente e codificar as respostas contra
a lista de achados deste roteiro, para que a fase qualitativa feche o
ciclo: cada mecanismo declarado nas entrevistas deve ser reexaminável no
painel.

\end{document}
